\documentclass[10pt]{article}
\usepackage{resume}
\usepackage{xeCJK}
\xeCJKDeclareCharClass{Default}{"2013}
\setCJKmainfont{FandolSong-Regular.otf}[BoldFont=FandolHei-Bold.otf,Script=Default]
\renewcommand{\rolestyle}{\upshape}
\renewcommand{\rolesep}{｜}
\hypersetup{pdftitle={李哲骧 - AI 工程岗位简历},pdflang={zh-CN}}
\begin{document}
\chineseheader

\section{教育背景}
\entry{加州大学洛杉矶分校（UCLA）}{预计 2028.06 毕业}
计算机科学硕士在读
\entry{南加州大学（USC）}{2026.05 毕业}
计算机科学学士（Magna Cum Laude 荣誉毕业）；GPA：3.76/4.00；Viterbi Dean's List

\section{技术能力}
\textbf{编程语言：}Python、C/C++、JavaScript/TypeScript、SQL、Java、Swift\\
\textbf{AI 开发：}PyTorch、Triton、大模型 API、智能体工作流、RAG（BGE 向量、ChromaDB）、多模态推理\\
\textbf{应用开发：}FastAPI、asyncio、SQLAlchemy、MySQL、PostgreSQL（Supabase）、Next.js、Docker、Linux、Git

\section{工程与科研经历}
\role{UCLA Health Intelligence Lab（HAIL）}{研究生研究员}{}{2026.08 – 至今}
\begin{itemize}
\item 在 Yuzhe Yang 教授指导下研究多模态血糖预测：已将 49 个数据集的 CGM 数据整理为统一格式，纳入饮食、胰岛素与可穿戴设备信息；正在搭建按患者划分的基准，用于比较预测模型。
\end{itemize}
\role{未来求索 / PMAOS}{AI 工程实习生}{}{2026.07 – 2026.08}
\begin{itemize}
\item 设计短视频推荐算法，将用户与视频双向兴趣传播、多时间尺度曝光预算和疲劳控制结合，用于管理重复推荐与不同兴趣的曝光分配。
\item 开发支持 11 种语言的音视频情绪识别，融合多模态证据，在模态冲突时暂不判断。
\end{itemize}
\role{Terlingua Threads}{全栈工程师}{远程}{2026.01 – 2026.05}
\begin{itemize}
\item 开发 Shopify 集成页面和 Google Maps 功能，完善移动端布局与无障碍访问。
\end{itemize}
\role{Piezolab Medical}{AI 工程实习生}{加州 Lake Forest}{2025.06 – 2025.08}
\begin{itemize}
\item 搭建大模型辅助流程，对 5,000 余份医疗事故资料进行导入、清洗与聚类，支持医疗质量分析；为血管护理研究开发智能体工作流程。
\end{itemize}

\section{代表 AI 项目}
\role{AI Doctor WeChat}{Python、FastAPI、MySQL、RAG}{}{2026}
\begin{itemize}
\item 开发并在微信云托管部署 FastAPI 后端，支撑诊所小程序与 Web 工作台的临床教学和患者信息采集。
\item 设计五类角色权限，大模型整理的病历须经医生确认后才对患者可见；实现 RAG（BGE 向量、ChromaDB、关键词兜底）和限流的 asyncio 大模型调用队列。
\end{itemize}
\role{Fridge Manager}{Next.js、TypeScript、Supabase、大模型与视觉模型}{}{2026}
\begin{itemize}
\item 开发并上线 AI 厨房助手 \link{https://howtoeat.food}{howtoeat.food}，根据库存、剩菜、偏好、厨具与烹饪时间生成菜单和采购清单。
\item 实现拍照识别库存与大模型菜谱导入（用户修正、管理员审核、去重）；搭建多模态视频流程（字幕、ASR/OCR、关键帧视觉），把烹饪视频整理成结构化菜谱。
\end{itemize}
\role{投资分析 Agent Skill}{Python、Claude Code}{}{2026.07}
\begin{itemize}
\item 编写智能体技能，根据公司名称或商业计划书调研并起草有证据支撑的投资报告；用 Python 质检框架检查每一节的来源、深度和完整性。
\end{itemize}

\section{代表算法研究}
\begin{itemize}
\item \textbf{Efficient Attention via Pre-Scoring}（第一作者，审稿中）：提出在近似注意力前筛选关键 key，用 Triton 实现，并与 HyperAttention 结合在 131k token 长上下文上评估。\link{https://arxiv.org/abs/2505.11040}{arXiv:2505.11040}
\item \textbf{MedTokenBudget}（独立作者，审稿中）：病灶监督的 token 路由，使皮肤镜图像分类在固定 token 预算下保留诊断区域，在 ISIC 2019 上评估。
\end{itemize}
\end{document}
